\chapter{Approximation theory}
This is the auxiliary, supplement note for now about polynomial structures and thereof, coming from chapter~\ref{chapter:double_descent_prob} in the problem section. So far, we will go all in, for whatever knowledge are there that can be known of polynomial. 
\section{Fundamental knowledge}

There are a few ways to express the structure of polynomial. Unlike the usual exploratory and in one way or another, newer foundational approach, we approach polynomial using the existing infrastructures, such is to say we assume and abide of many assumptions and structures thereof of the values field it is on. Let us begin with the basic. 

The first case of any given functional structures, are that they act on the structure of a given set, for example in this case, usually it is $\mathbb{R}$. Thereby, polynomial is the structure on which relations are made between equality operation $=$, perhaps inequality $<,>,\geq,\leq,\not\leq,\not\geq$, and so on, and with variables and coefficients such is to form the expression in which different transformation are indicted upon. The equality, as stated, is extremely necessity as a tool of abstraction, though in this case it works as abstraction and not evaluation thereof, and thus can be considered an abuse of notation, written such way to streamline understanding of which lies trivial in most case (nevertheless, we have to take into account such notion itself). So, what can be considered a polynomial? $3t^{3}-7t^{2}+4t-1$ and $8t^{6}-t^{5}+\sqrt{2}t^{2}+\sqrt{-1}t-1$ are all polynomials. So are $(t^{2}-1)/(t-1)$ for $t\neq 1$. What can be the definition? 

\begin{definition}[Polynomial, preliminary definition]
    A function of a single variable $t$ is a polynomial on its domain if we can put it in the form 
    \begin{equation}
        a_{n}t^{n} + a_{n-1}t^{n-1} + \dots + a_{1}t + a_{0}t^{0}
    \end{equation}
    where $\{a_{i}\}$ are constants, and of a singular variable $t$.  
\end{definition}
We note here that we are defining this for \textit{single-variable} polynomial only. We notice earlier in our statement about polynomial structure, there exists this line of argument in terms of single or multi-variable, as such, we can extend upon. Using the equality sign, we get the assignment formal computation for a polynomial of $m$ variables as \begin{equation}
    p(n,m) = p_{n}(x_{1},x_{2},\dots,x_{m})
\end{equation}
This expression describe the degree $n$ and the number of variable $m$. Let us restrict to $m=1$ case of single-variable. Then, the expression can be naturally defined as 
\begin{equation}
    p_{nm}(x) = p_{n}(x) = \sum_{n\leq i \leq 0} a_{i}x_{i} = a_{n}x^{n} + a_{n-1}x^{n-1} + \dots + a_{1}x + a_{0}x^{0}
\end{equation}
and so on. The notation is shifted for convenience, and this is the standard notation in the single-variable case. Interestingly, the definition is very naturally defined for single-variable, we will get to this later on with the multi-variable approach. Any expression that can be converted, or transformed analytically to similar form is said to be a polynomial. Hence, for something as $(t^{2}-1)/(t-1)$, we can convert it to $t+1$, though there exists a fundamental difference that the expression in the conversion must satisfy $x\neq 1$ for completeness, else they would not be equal. This is the fundamental fact about equality and elementary expression transformation. One can shift from $(t^{2}-1)/(t-1)$ to its reduced form by various axioms and rules to $t+1$, but the structure itself encodes $t\neq 1$ into it. On the other hand, if you go from $t+1$ to $(t^{2}-1)/(t-1)$, such constraint is inverted, and hence you are reaching for semi-equivalent class of expression instead. This topic goes much deeper, so we will leave it here for purposes. Though indeed, from the viewpoint of a transformation, this is a transformation with losing information. 

For the polynomial, the numbers $a_{i}$ for $0\leq i \leq n$ are called \textit{coefficient}\index{polynomial coefficient}. $a_{n}$ is called the \textbf{leading coefficient}, and $a_{n}t^{n}$ the leading term of the polynomial. When the leading coefficient $a_{n}=1$, we call the polynomial to be \textit{monic}. Monic polynomial itself is fairly important, as for many of its convenience conventions exist. 

The non-negative integer $n$ is the degree of the polynomial, we write such as $\deg{p}=n$. We also recall from elementary mathematics, $$a^{n}=\underbrace{aaa\dots a}_{n\text{ times}}$$ 
A \textit{constant} polynomial has but a single term, $a_{0}$, and a nonzero constant polynomial has degree 0, though per convention, the zero polynomial has degree $-\infty$. A \textit{zero} of a polynomial $p(x)$ is any number $r$ for which $p(r)=0$. Such is then called the \textit{root} \index{polynomial root} of the polynomial, or a solution of the equation $p(r)=0$. There are many solutions, and cases that require such root, for example, for a polynomial to pass through a point $y$, so that, $p(x)=y$, then $p(x)-y=0$ thus create an extended polynomial form of $p'(x)=p(x)-y$. Most problem can be reduced to such form. In some cases, the solution does not exist within the set of real numbers itself, but rather in the set of \textit{complex number} $\mathbb{C}$ for all $a+bi \in \mathbb{C}$, such that $i^{2}=1$. This is somewhat particularly troubling, as even though encoding exists in mathematics per our view, emerging from solving realized and principle-ly real system to complex domain is still an open question. 

In operating with polynomials, we can treat them in somewhat a same way as scaling the operator into acting with polynomials. 

\begin{theorem}[Polynomial operations]
    Extend the elementary operations to polynomial object $p,q,\dots$. Let $p_{n}(x)$ and $q_{m}(x)$ be any two polynomials of degree $n$ and $m$. Then. 
    \begin{enumerate}[leftmargin = 2cm]
        \item The sum $(p+q)(x)$ is evaluated to $(a_{0}+b_{0})+(a_{1}+b_{1})x+(a_{2}+b_{2})x^{2}+\dots$ if $m=n$, the equation matches. 
        \item The difference $(p-q)(x)$ is evaluated to $(a_{0}-b_{0})+(a_{1}-a_{1})x+\dots$. 
        \item The product of a constant and a polynomial is defined as $cp(x)=ca_{0}+ca_{1}x+\dots$. 
        \item The product of two polynomials is defined, for $(pq)(x)$, as followed: 
        \begin{equation}
            \begin{split}
                (pq)(x) &= a_{0}b_{0} + (a_0 b_{1} + a_{1}b_{0})x + (a_0 b_{2} + a_{1}b_{1}+ a_{2}b_{0})x^{2} + \dots\\
                & \quad   + (a_{0}b_{r}+a_{1}b_{r-1}+\dots+ a_{i}b_{r-i}+\dots+ a_{r}b_{0})t^{r} +\dots + (a_{n}b_{m})x^{m+n}
            \end{split}
        \end{equation}
        \item Composition between $p$ and $q$ is $(p\circ q )(x)= p(q(x))$. This definition instruct us to replace each occurrence of $x$ in the expression to $q(x)$. The order definitely matters. 
    \end{enumerate}
\end{theorem}

\subsection{Polynomial of several variables}

If $x$ and $y$ are the roots of a quadratic equation $at^{2}+bt+c=0$, then $-b/a=a+y$ and $c/a=xy$. The expressions $x+y$ and $xy$ are example of polynomial-like expression with two variables $x$ and $y$. In general, a function $f(x,y)$ is a \textit{multivariate polynomial} in $x,y$ if and only if it can be represented as a finite sum of terms of the form $cx^{k}y^{m}$, where $c$ is a coefficient and $k$ an $m$ are nonnegative integers. The number $k+m$ is called the degree of the term, and the degree of the polynomial $f(x,y)$ is equal to the highest degree of its own terms. Polynomial of several variables can be addled, subtracted and multiplied in a way analogous to polynomial of a single variable. There are two classes of multivariate, and in specific, two-variable polynomial that we want to concern of: 
\begin{enumerate}[label=\roman*]
    \item Symmetric polynomials $f(x,y)$ in which $f(x,y)=f(y,x)$. 
    \item Homogeneous polynomials in which all the terms are of the same degree. 
\end{enumerate} 
Similar definition can be made for functions of three variables, say $x,y,z$. We can see the shift between defining and structuring polynomial of several variables here. By default, the concept of polynomial is applicable mainly for single variable. When it comes to this expansion, it is inherently problematic because of its own criteria supporting the definition of polynomial itself. For example, what is the degree for this new multivariate polynomial, and so on. In fact, for any given degree, there can be infinitely many constructions that give such such degree, with respect to a variable explicit degree. 
\subsection{Complex numbers}
Many solutions exist purely on the basis of complex number. Originally, complex numbers are regarded as a quick guesswork that make it possible so that certain real-valued equation works toward solution, for example, $x^{2}+2x+3$. However, it faced heavy criticism against the physical reality of certain quantification by context, such as counting, measuring distance relations, and so on. However, if we encode the context in a different way, and provide it a different representation, sometimes bringing it to simply be an abstract mathematical language, such work is fairly interesting, and sometimes particularly intriguing. Toward such stance, it is wise for us to take the abstract, and continue on discussing about complex number. 

Let $z=x+iy$ denotes a \textit{complex number}\index{complex number} with $x$ and $y$ real. The real part, denoted $\mathrm{Re}{(z)}$ is the number $x$, while the imaginary part, $\mathrm{Im}{(z)}$ is the number $y$. Such is then to say that complex numbers are formed using the combination of two real set, so $\mathbb{C}\equiv \mathbb{R}\times\mathbb{R}$ with added structure as the imaginary unit $i$, in which $i=\sqrt{-1}$. The meaning of $i$ is then implicit and is often waived for `simplicity', however, looking from certain perspective, $i$ represents certain kind of rotation. The \textit{complex conjugate}, $\bar{z}$ of $z$ is $x-yi$, while the \textit{modulus} or \textit{absolute value} of $z$, denoted by $|z|$, is defined as $\sqrt{x^{2}+y^{2}}$. This is the distance from the origin to the point representing $z$. The argument of $z$ is the angle between the real axis and the line joining $0$ and $z$, measured in the counterclockwise direction. It is denoted by $\mathrm{arg}{(z)}$, and generally, we assign $\mathrm{arg}{(z)}$ to be a value between $0$ and $2\pi$. From such, we then also have the \textit{polar decomposition}. Let $r=|z|, \theta = \mathrm{arg}{(z)}$, then we have 
\begin{equation}
    z= r(\cos{(\theta)}+ i\sin{(\theta)})
\end{equation}
Furthermore, via certain mathematical eldritch ritual of which connecting multiple interpretations and structures together, Euler came up with the law for complex representation, 
\begin{equation}
    e^{it} = \cos{(t)} + i\sin{(t)}, \quad t\in \mathbb{R}
\end{equation}
One then can convert one form to the other by using Euler's formula and vice versa to polar decomposition:
\begin{equation}
    z= x+ iy \Leftrightarrow z = \sqrt{x^{2}+y^{2}} \exp{i\left(tan^{-1}{\left(\frac{y}{x}\right)}\right)} = \rho e^{i\theta}
\end{equation}
Usually, we restrict $0\leq \theta \leq 2\pi$ or $-\pi \leq \theta \leq \pi$. Now let us come back to the problem of the equation, for such, $x^{2}+2x+3$. Now, what did we say when we examine it previously? We choose the set or field $\mathbb{R}$ as the solution space, and only in such space, the equation has \textit{no solution}. There is no way to find such root from such polynomial in real space. However, if we extend the notion of the number system in working to complex number, which encodes rotation by $i$ axis, then intrinsically, we have some satisfactory solution in play. So now, the expression is more like 
\begin{equation}
    (a+bi)^{2} + 2(a+bi)+ 3 = 0
\end{equation}
Of which now rotation is possible. So one would just "rotate" the function for intersection, and find out that $-1 \pm i\sqrt{2}$ is the answer on the complex space. Thereby, solving for polynomial, or as any type of functions regardless of such, will have solutions depends on what space you confine the object in question in. 

\section{Algebraic polynomial}

Let $K$ be, then, an arbitrary field, and consider polynomials whose coefficients are element of $K$. The typical polynomial of this kind is as usual, with each of the coefficients $a_{i}$ being an element of $K$. For economy of notation, we might as well use $f(x)$ to denote polynomial and add details subscripts for additional insight. The set of all polynomials with coefficient in $K$ will be written as $K[x]$. The element of $K$ themselves can be regarded as polynomials, namely as the constant term polynomial, thus calling themselves the \textit{constant polynomial}. All the operations possible work the same here. A polynomial $a_{n}x^{n}+\dots+a_{0}$ i $K[x]$ with $a_{n}\neq 0$ is said to have degree $n$. We have some certain facts about the $K[x]$ ring. 

\begin{theorem}[Integral domain]\index{integral domain}
    Let $K$ be a field, then we have:
    \begin{enumerate}[leftmargin=2cm,topsep=1pt,itemsep=1pt]
        \item $K$ is an integral domain. 
        \item For two polynomial $f(x)$ of degree $m$ and $g(x)$ or degree $n$ in $K[x]$, the product $f(x)g(x)$ has degree $m+n$. 
        \item $K[x]$ is an integral domain. 
        \item A polynomial $f(x)$ in $K[x]$ is a unit in $K[x]$ if and only if $f(x)$ has degree $0$. 
    \end{enumerate}
\end{theorem}

In the language of ring, of which polynomial class is held upon, a ring $R$ is said to contain \textit{zero-divisor} if there exists one or more nonzero elements that when multiplied together give the result $0$. If not, then it is then called the \textit{integral domain}. In such sense, we can think of integral as meaning that in an integral domain, the multiplication operation maintains its integrity, without suddenly collapsing to the respective identity element. 

\begin{theorem}[Factoring a polynomial]
    Let $K$ be a field. A polynomial $f(x)$ of positive degree in $K[x]$ either is irreducible in $K[x]$ or factors as a product of irreducible polynomials in $K[x]$. 
\end{theorem}

We have certain result using factorization of a polynomial toward the root of such. 
\begin{theorem}[Polynomial remainder theorem]
    Let $K$ be a filed. If $a(x)$ and $b(x)$ are nonzero polynomials in $K[x]$, then there exist unique polynomial $q(x)$ and $r(x)$ with $r(x)$ having lower degree than $a(x)$, such that 
    \begin{equation}
        b(x) = a(x)q(x) + r(x)
    \end{equation}
    We call $q(x)$ the quotient polynomial, and $r(x)$ the remainder polynomial obtained on dividing $b(x)$ by $a(x)$. 
\end{theorem}

\subsection{Minimal polynomials}

Suppose that $K$ is a field and $f(x)$ a polynomial in $K[x]$. Then polynomial $f(x)$ may have no roots in $K$ yet have roots in a larger field $L$. In general, the study of polynomial $f(x)$ with coefficients in field $K$ naturally leads to the study of larger fields that contain $K$, these larger fields being the source of missing roots of the polynomials. 

A \textit{field extension} is a pair of fields $(K,L)$ such that $K$ lies inside $L$ and such that the rules of addition and multiplication for $K$ and $L$ are compatible. That is, for two elements of $K$, their sum and product as members of $K$ coincide with their sum and product as members of $L$. We often use the notation $K\subseteq L$ to describe a field extension consisting of the field $K$ and $L$ with $K$ inside $L$. We then call $L$ a \textit{field extension} of $K$. Alternatively, one may start with a field extension $K\subseteq L$, in which case $s\in L$ lead in a natural way to polynomial in $K[x]$, the polynomial that have these elements as root. 

For a field extension $K\subseteq L$, to study such extension in natural form, we come up with the concept of \textit{algebraic} over $K$. An element $\gamma$ of $L$ is algebraic over $K$ if there is a nonzero polynomial $f(x)$ in $K[x]$ such that $f(\gamma)=0$. Thus, $\gamma$ is algebraic if there are elements $a_0, a_{1},\dots,a_{n}$ of $K$ with $n>0$ and $a_{n}\neq 0$ such that 
\begin{equation}
    a_{n}\gamma^{n} + a_{n-1}\gamma^{n-1} + \dots + a_{1}\gamma^{1} + a_{0} = 0
\end{equation}
Let $K\subseteq L$ be a field extension, and suppose that $\gamma$ is an element of $L$ and that $f(x)$ is a polynomial in $K[x]$ such that $\gamma$ is a root. If $g(x)$ is an arbitrary polynomial in $K[x]$, the product $f(x)g(x)$ will also have $\gamma$ as a root, $g(\gamma)f(\gamma)=0g(\gamma)=0$. Thus there is a large collection of polynomials of $K[x]$ that have $\gamma$ as root, not just the given one $f(x)$. Despite such, it is possible to give a succinct description of the polynomials in it. All the polynomial in $K[x]$ with $\gamma$ as a root are multiples of a polynomial of least degree with $\gamma$ as a root, which is then called the \textit{minimal polynomial} of $\gamma$ over $K$. \index{minimal polynomial}

Two results in such analysis are useful, with one is a refinement of the other. 
\begin{theorem}
    Suppose $K\subseteq L$ is a field extension and $\gamma$ is an element of $L$ that is algebraic over $K$ of degree $n$. Let $f(x)$ be an $n$-degree polynomial in $K[x]$ such that $f(\gamma)=0$ and let $g(x)$ be an arbitrary polynomial in $K[x]$ such that $g(\gamma)=0$. Then $f(x)$ divides $g(x)$. 
\end{theorem}

\begin{theorem}
    Suppose $K\subseteq L$ is a field extension and $\gamma$ is an element of $L$ that is algebraic over $K$ of degree $n$. Then, 
    \begin{enumerate}[leftmargin = 2cm, topsep = 1pt, itemsep = 1pt]
        \item There is a unique monic polynomial $p_{\gamma}(x)$ of degree $n$ in $K[x]$ having $\gamma$ as a root. 
        \item The polynomial of $K[x]$ that have $\gamma$ as a root are precisely the polynomials divisible by $p_{\gamma}(x)$. 
        \item The polynomial $p_{\gamma}(x)$ is irreducible in $K[x]$. 
    \end{enumerate}
\end{theorem}
\section{Approximation theory}

We would like to investigate more, after the structural analysis, into the role of analysis on the function of polynomial. Polynomial on itself is nothing that special, in an analytic sense, though they are practically fairly useful. Construction of a polynomial then is one of the main focus in analysis, as to inspect what kind of polynomial can there be aside from the ring structure $K[x]$ itself. This prompted two kinds of direction in a polynomial study: Can polynomial be used to examine, extrapolate and capture different concepts or functions of arbitrary space; and what can be all type of added structure polynomial construction there are that is constructed from such algebraic structure? One of those questions can be answered on the basis of \textit{function approximation}, by Weierstrass Approximation Theorem. 

At the age of 70, Karl Weierstrass (Karl Theodor Wilhelm Weierstrass --- 1815-1897) published the result in mathematical analysis of which is the Weierstrass approximation theorem, which elementarily states that every continuous function defined on closed interval $[a,b]$ can be uniformly approximated closely as desired by a polynomial function. The theorem is often called the Stone-Weierstrass theorem for the contribution of Marshall H. Stone on generalization of the theorem itself, on the compact Hausdorff space and with arbitrary continuous function class. Nevertheless, we should be acquainted of either of its form, and the implication thereof. There are several forms of the theorem, so we will have to take the long form to analyse such. 

Of course, it is not limited to simply polynomial, but to a wide range of different functions and thereof. In fact, approximation can sometimes be attributed to P. L. Chebyshev, who while working in 1853 with the problem of \textit{linkage}, devices which translate the linear motion of a steam engine into circular motion of a wheel, considered the problem of approximation:
\begin{quote}
    Given a continuous function $f$ defined on a closed interval $[a,b]$ and a positive integer $n$, can we `represent' $f$ by a polynomial $p(x)=\sum_{k=0}^{n}a_{k}x^{k}$ of degree at most $n$, in such a way that the maximum error at any point $x$ in $[a,b]$ is controlled? In particular, is it possible to construct $p$ so that the error $\max_{a\leq x\leq b}|f(x)-p(x)|$ be minimized? 
\end{quote}
This particular problem is the fairly impressive approximation problem at first, of which raises particular problems and else that we would soon see in analysis: 
\begin{enumerate}[topsep=1pt,itemsep=1pt]
    \item Why should such a polynomial even exists? 
    \item If it does, can we hope to construct it? 
    \item If it exists, is it also unique? 
    \item What happens if we change the measure of error to something else, for example, the average mean square $\int_{a}^{b}|f(x)|-p(x)|^{2}\:dx$?
    \item How much can the minimization be achieved? 
\end{enumerate}

Those problems are fundamental of the approximation theory. 

